\documentclass{article}
\usepackage[utf8]{inputenc}
\usepackage{amsmath,amssymb}
\usepackage[most]{tcolorbox}
\usepackage{geometry}
\geometry{margin=2.5cm}

\newcommand{\step}[1]{\par\noindent\hangindent=3.5em\hangafter=1%
  \makebox[3.5em][l]{\textbf{#1.}}}
\newcommand{\steptag}[1]{\hfill{\small\textsf{[#1]}}}

\newtcolorbox{proofbox}[1]{
  colback=gray!5, colframe=gray!60,
  fonttitle=\bfseries, title={#1},
  breakable, enhanced,
  left=4pt, right=4pt, top=4pt, bottom=4pt,
  before skip=10pt, after skip=10pt
}

\begin{document}

\begin{proofbox}{Uniform unitary graph control: there are universal C\_ch >...}
\step{1} Uniform unitary graph control: there are universal C\_ch >= 1, kappa\_ch in (0, 1/2] such that for every finite-dimensional exact-unit epsilon\_r-C*-algebra, every delta > 0 with C\_ch*(epsilon\_r + delta) <= kappa\_ch, every V in calUbar\_delta, and every A\textasciicircum{}par in B\textasciicircum{}\{icalH\}\_\{2delta\}(0), there is a unique g\_V(A\textasciicircum{}par) in B\textasciicircum{}\{calH\}\_\{2delta\}(0) with f\_V(A\textasciicircum{}par + g\_V(A\textasciicircum{}par)) = 0, and the corresponding V bold-dot (J + A\textasciicircum{}par + g\_V(A\textasciicircum{}par)) lies in calU, where f\_V(A) = (1/2)*(((J + A\textasciicircum{}dagger) bold-dot V\textasciicircum{}dagger) bold-dot (V bold-dot (J + A)) - J); moreover ||g\_V(A\textasciicircum{}par) + (1/2)*(V\textasciicircum{}dagger bold-dot V - J)|| <= C\_ch*(epsilon\_r*delta + delta\textasciicircum{}2), ||Dg\_V(A\textasciicircum{}par)|| <= C\_ch*(epsilon\_r + delta), and ||D\_\{A\textasciicircum{}perp\} f\_V(A\textasciicircum{}par + g\_V(A\textasciicircum{}par)) - I\_\{calH\}|| <= C\_ch*(epsilon\_r + delta) < 1; the resulting C\textasciicircum{}1 graph charts cover calU. \steptag{validated} \\[6pt]
\step{1.1} Uniform multiplier control: write e=epsilon\_r and s=e+delta. If 0<=e,delta and s<=1/512, V lies in calUbar\_delta, and q=V\textasciicircum{}dagger bold-dot V-J, then ||V||<=1+2s, L\_V is invertible, ||L\_V\textasciicircum{}(-1)||<=2, and ||L\_V\textasciicircum{}(-1)-L\_\{V\textasciicircum{}dagger\}||<=8s. Moreover, whenever A has ||A||<=4delta, ||L\_V\textasciicircum{}(-1)(L\_\{V bold-dot (J+A)\}-L\_V)||<=16s<1, so L\_\{V bold-dot (J+A)\} is invertible and V bold-dot (J+A) has a right inverse. \steptag{validated} \\[6pt]
\step{1.1.1} Base multiplier estimate from the exact-unit epsilon\_r-C*-axioms. Let Y be a right inverse of V, so V bold-dot Y=J. Since ||q||<=2delta and the C*-lower bound gives (1-e)||V||\textasciicircum{}2<=||V\textasciicircum{}dagger bold-dot V||<=1+2delta, s<=1/512 implies ||V||<=1+2s. Approximate associativity and exact unitality give ||V\textasciicircum{}dagger-Y||<=e||V||\textasciicircum{}2||Y||+(1+e)||q||||Y||<=4s||Y||, hence ||Y||<=||V||/(1-4s)<=2. Thus ||L\_V L\_Y-I||<=e||V||||Y||<1, so L\_V has a right inverse. Also ||L\_\{V\textasciicircum{}dagger\}L\_V-I||<=e||V||\textasciicircum{}2+(1+e)||q||<=4s<1; therefore (L\_\{V\textasciicircum{}dagger\}L\_V)\textasciicircum{}(-1)L\_\{V\textasciicircum{}dagger\} is a left inverse. The left and right inverses coincide, and the Neumann bounds yield ||L\_V\textasciicircum{}(-1)||<=2 and ||L\_V\textasciicircum{}(-1)-L\_\{V\textasciicircum{}dagger\}||<=[4s/(1-4s)](1+e)||V||<=8s. \steptag{validated} \\[4pt]
\step{1.1.2} Perturbation to the chart point. For ||A||<=4delta put X=V bold-dot (J+A)=V+V bold-dot A. The product-norm axiom and node 1.1.1 give ||L\_X-L\_V||<= (1+e)||X-V||<=4(1+e)\textasciicircum{}2(1+2s)delta<=8delta, whence ||L\_V\textasciicircum{}(-1)(L\_X-L\_V)||<=16delta<=16s<1. The Neumann lemma makes L\_X=L\_V[I+L\_V\textasciicircum{}(-1)(L\_X-L\_V)] invertible. Solving L\_X(Y\_X)=J then gives X bold-dot Y\_X=J, so X has a right inverse. \steptag{validated} \\[6pt]
\step{1.1.2.1} Dependency-gated perturbation argument. Require validated node 1.1.1. Under the shared hypotheses e,delta >= 0, s=e+delta <= 1/512, V in calUbar\_delta and ||A||<=4delta, node 1.1.1 supplies ||V||<=1+2s, invertibility of L\_V, and ||L\_V\textasciicircum{}(-1)||<=2. Exact unitality and bilinearity give X=V bold-dot(J+A)=V+V bold-dot A. Hence the product-norm axiom gives ||L\_X-L\_V||<= (1+e)||X-V||<= (1+e)\textasciicircum{}2||V||||A||<=4(1+e)\textasciicircum{}2(1+2s)delta<=8delta, where the last inequality follows from e<=s<=1/512. Therefore B=L\_V\textasciicircum{}(-1)(L\_X-L\_V) satisfies ||B||<=16delta<=16s<=1/32<1. The Neumann series makes I+B invertible; the exact operator identity L\_X=L\_V(I+B) then makes L\_X invertible. Setting Y\_X=L\_X\textasciicircum{}(-1)(J) gives X bold-dot Y\_X=L\_X(Y\_X)=J, so X has a right inverse. \steptag{validated} \\[4pt]
\step{1.2} Explicit graph-equation estimates: under the hypotheses of node 1.1, define q=V\textasciicircum{}dagger bold-dot V-J and F(P,H)=f\_V(P+H) on anti-Hermitian P and Hermitian H with ||P||,||H||<2delta. Then F is a C\textasciicircum{}1 map of finite-dimensional real spaces into calH and, with the explicit universal K=64, ||F(P,H)-(q/2+H)||<=K*(epsilon\_r*delta+delta\textasciicircum{}2), ||D\_H F(P,H)-I\_calH||<=K*(epsilon\_r+delta), and ||D\_P F(P,H)||<=K*(epsilon\_r+delta). \steptag{validated} \\[6pt]
\step{1.2.1} Zero-order expansion with no hidden big-O term. Put A=P+H, so ||A||<4delta and A+A\textasciicircum{}dagger=2H. Bilinearity gives 2F(P,H)=V\textasciicircum{}dagger bold-dot V-J+V\textasciicircum{}dagger bold-dot(V bold-dot A)+(A\textasciicircum{}dagger bold-dot V\textasciicircum{}dagger) bold-dot V+(A\textasciicircum{}dagger bold-dot V\textasciicircum{}dagger) bold-dot(V bold-dot A). Reassociate only the two linear terms and write V\textasciicircum{}dagger bold-dot V=J+q; then 2F=q+2H+R, where R=q bold-dot A+A\textasciicircum{}dagger bold-dot q+r\_1+r\_2+(A\textasciicircum{}dagger bold-dot V\textasciicircum{}dagger) bold-dot(V bold-dot A), ||r\_1||,||r\_2||<=e||V||\textasciicircum{}2||A||. Using ||q||<=2delta, ||V||<=1+2s from node 1.1.1, the product bound, s<=1/512, and ||A||<4delta gives ||R||<=128*(e*delta+delta\textasciicircum{}2). Hence ||F(P,H)-(q/2+H)||<=64*(e*delta+delta\textasciicircum{}2). \steptag{validated} \\[4pt]
\step{1.2.2} Derivative expansion. For arbitrary Z, exact bilinearity and the exact reversal law give 2DF(P,H)[Z]=((Z\textasciicircum{}dagger bold-dot V\textasciicircum{}dagger) bold-dot(V bold-dot(J+A)))+(((J+A\textasciicircum{}dagger) bold-dot V\textasciicircum{}dagger) bold-dot(V bold-dot Z)). Reassociate the terms independent of A through V\textasciicircum{}dagger bold-dot V=J+q; their errors are bounded by e||V||\textasciicircum{}2||Z|| and their q-terms by 2(1+e)delta||Z||. Bound every term containing A directly using ||A||<4delta, ||V||<=1+2s, and the product norm. The two displayed summands then differ from Z\textasciicircum{}dagger and Z, respectively, by at most 32s||Z|| each, so ||DF(P,H)[Z]-(Z+Z\textasciicircum{}dagger)/2||<=32s||Z||<=64s||Z||. Restriction to Hermitian Z gives ||D\_HF-I||<=64s, while restriction to anti-Hermitian Z gives ||D\_PF||<=64s. Finally, with X=V bold-dot(J+A), exact product reversal identifies 2F=X\textasciicircum{}dagger bold-dot X-J, which is Hermitian; the displayed degree-two real polynomial is therefore C\textasciicircum{}1 from icalH times calH to calH. \steptag{validated} \\[4pt]
\step{1.3} Quantitative graph construction: if additionally c=64*(epsilon\_r+delta)<=1/8, then for every P in B\textasciicircum{}\{icalH\}\_\{2delta\}(0) there is a unique g\_V(P) in B\textasciicircum{}\{calH\}\_\{2delta\}(0) with F(P,g\_V(P))=0. These values form a C\textasciicircum{}1 map and satisfy ||g\_V(P)+q/2||<=64*(epsilon\_r*delta+delta\textasciicircum{}2), ||Dg\_V(P)||<=128*(epsilon\_r+delta), and ||D\_H F(P,g\_V(P))-I\_calH||<=64*(epsilon\_r+delta)<1. \steptag{validated} \\[6pt]
\step{1.3.1} Existence and uniqueness by the exact external lem-stage1-quantitative-inverse-function. Fix P and apply that lemma to h↦F(P,h) on B\textasciicircum{}\{calH\}\_\{2delta\}(0), with the Banach-space isomorphism I\_calH, center 0, radius 2delta, and derivative constant c=64s. Node 1.2 gives ||D\_HF-I||<=c<1, so the map is injective and its image contains F(P,0)+B\textasciicircum{}\{calH\}\_\{2delta(1-c)\}(0). Also node 1.2.1 and ||q||<=2delta give ||F(P,0)||<=delta+64(e delta+delta\textasciicircum{}2)=delta(1+c)<2delta(1-c), since c<=1/8. Hence 0 belongs to that image and determines a unique g\_V(P) in the stated open ball. \steptag{validated} \\[4pt]
\step{1.3.2} Quantitative and C\textasciicircum{}1 dependence. Substituting H=g\_V(P) in node 1.2.1 and using F(P,g\_V(P))=0 gives ||g\_V(P)+q/2||<=64(e delta+delta\textasciicircum{}2); node 1.2.2 directly gives the normal derivative estimate. For P\_1,P\_2, the lower Lipschitz inequality in lem-stage1-quantitative-inverse-function for h↦F(P\_1,h), followed by integration of ||D\_PF||<=c along the segment from P\_1 to P\_2, gives (1-c)||g\_V(P\_1)-g\_V(P\_2)||<=c||P\_1-P\_2||. Thus g\_V is locally Lipschitz. Expanding the polynomial F at (P,g\_V(P)) and using that Lipschitz bound shows that its Frechet derivative is Dg\_V(P)=-(D\_HF(P,g\_V(P)))\textasciicircum{}(-1)D\_PF(P,g\_V(P)). The Neumann bound is ||(D\_HF)\textasciicircum{}(-1)||<=1/(1-c), so ||Dg\_V(P)||<=c/(1-c)<=2c=128s. Continuity of DF, of g\_V, and of inversion on invertible operators makes Dg\_V continuous; hence g\_V is C\textasciicircum{}1. \steptag{validated} \\[4pt]
\step{1.4} Unitary graph, covering, and constants closure: under nodes 1.1-1.3, X\_V(P)=V bold-dot(J+P+g\_V(P)) satisfies X\_V(P)\textasciicircum{}dagger bold-dot X\_V(P)=J and has a right inverse, hence belongs to calU; the C\textasciicircum{}1 injective graphs describe the local unitary locus and cover calU. Choosing the universal constants C\_ch=1024 and kappa\_ch=1/4 makes the root guard imply every smallness hypothesis in nodes 1.1-1.3, and their bounds strengthen to all three advertised C\_ch estimates with strict normal-derivative bound <1. \steptag{validated} \\[6pt]
\step{1.4.1} Exact graph, cover, and constant closure. For A=P+g\_V(P), exact product reversal gives X\_V(P)\textasciicircum{}dagger bold-dot X\_V(P)-J=2F(P,g\_V(P))=0; since ||A||<4delta, node 1.1.2 makes L\_\{X\_V(P)\} invertible and supplies a right inverse, so X\_V(P) lies in calU by def-approximate-unitary-space. Since L\_V is invertible, equality of two graph points forces equality of P+g\_V(P), hence of their anti-Hermitian parts P; the inverse chart is phi\_V\textasciicircum{}parallel. Conversely, every unitary in the coordinate box has f\_V(A)=0 and the automatic right inverse, so node 1.3 identifies it with this graph. Every U in calU is covered by its self-centered chart because U lies in calUbar\_delta, f\_U(0)=0, and uniqueness gives g\_U(0)=0. Finally set C\_ch=1024 and kappa\_ch=1/4. The root guard gives s<=1/4096<1/512 and 64s<=1/64<1/8, so nodes 1.1-1.3 apply; their bounds 64(epsilon\_r delta+delta\textasciicircum{}2), 128s, and 64s are at most the corresponding C\_ch bounds, with 64s<=C\_ch s<=1/4<1. Thus the graph, cover, existence, uniqueness, membership, and all quantitative clauses of node 1 follow. \steptag{validated} \\[6pt]
\step{1.4.1.1} The epsilon\_r=0 boundary is included and valid. Node 1.1 assumes "0<=e,delta", i.e. 0<=e and 0<=delta, not 0<e; hence it and descendants 1.1.1-1.1.2 apply at e=0. Concretely there s=delta, exact associativity and exact unitality hold, and for a right inverse Y of V and q=V\textasciicircum{}dagger bold-dot V-J one has V\textasciicircum{}dagger=(V\textasciicircum{}dagger bold-dot V) bold-dot Y=Y+q bold-dot Y. Thus ||V\textasciicircum{}dagger-Y||<=||q||||Y||<=2delta||Y||; also L\_\{V\textasciicircum{}dagger\}L\_V=I+L\_q with ||L\_q||<=||q||<=2delta<1 under s<=1/4096. Together with the right inverse induced by Y this makes L\_V invertible, with the same (indeed stronger) bounds used in 1.1.2. The perturbation estimate there uses only delta<=s and 16s<1, so it also applies at e=0 and gives a right inverse for X\_V(P). Nodes 1.2 and 1.3 are stated under node 1.1 and likewise include e=0; their error terms simply specialize to O(delta\textasciicircum{}2) and O(delta). Therefore the root guard covers epsilon\_r=0 with no omitted case. \steptag{validated} \\[4pt]
\step{1.4.1.2} Dependency-gated graph and constants closure. Require validated nodes 1.1.1, 1.1.2.1, 1.3.1, and 1.3.2. Set C\_ch=1024 and kappa\_ch=1/4, and write e=epsilon\_r, s=e+delta. The root guard gives s<=1/4096<1/512 and c=64s<=1/64<1/8, so all four dependencies apply. For P in B\textasciicircum{}\{icalH\}\_\{2delta\}(0), nodes 1.3.1 and 1.3.2 give the unique H=g\_V(P) in B\textasciicircum{}\{calH\}\_\{2delta\}(0), its C\textasciicircum{}1 dependence, and the bounds ||H+q/2||<=64(e*delta+delta\textasciicircum{}2), ||Dg\_V(P)||<=128s, and ||D\_HF(P,H)-I\_calH||<=64s. Put A=P+H, so ||A||<4delta. Exact product reversal and F(P,H)=0 give X=V bold-dot(J+A) with X\textasciicircum{}dagger bold-dot X-J=2F(P,H)=0; node 1.1.2.1 gives a right inverse for X, hence X lies in calU by def-approximate-unitary-space. Node 1.1.1 gives invertibility of L\_V, and exact unitality gives phi\_V(X)=L\_V\textasciicircum{}(-1)(X-V)=A, so phi\_V\textasciicircum{}parallel(X)=P. Thus the graph map is injective. Conversely, if X=V bold-dot(J+A) is unitary with ||A\textasciicircum{}parallel||,||A\textasciicircum{}perp||<2delta, then exact product reversal gives F(A\textasciicircum{}parallel,A\textasciicircum{}perp)=0, and uniqueness in 1.3.1 forces A\textasciicircum{}perp=g\_V(A\textasciicircum{}parallel); therefore the graph is exactly the local unitary locus. For every U in calU, take V=U: then U lies in calUbar\_delta, F(0,0)=0, and uniqueness gives g\_U(0)=0, so U=X\_U(0), proving the cover. Finally 64<=C\_ch and 128<=C\_ch yield the first two advertised bounds, while ||D\_HF-I||<=64s<=C\_ch*s<=1/4<1 gives the third bound with strictness. This proves every clause of node 1 using only the listed validated dependencies. \steptag{validated} \\[4pt]
\step{1.4.2} Graph-chart identification and cover. Since L\_V is invertible, X\_V(P\_1)=X\_V(P\_2) implies P\_1+g\_V(P\_1)=P\_2+g\_V(P\_2); uniqueness of the Hermitian/anti-Hermitian decomposition gives P\_1=P\_2, so the C\textasciicircum{}1 parametrization is injective and its inverse is the continuous linear coordinate phi\_V\textasciicircum{}parallel. More generally, if X=V bold-dot(J+A) lies in the coordinate box ||A\textasciicircum{}parallel||,||A\textasciicircum{}perp||<2delta, then X is unitary exactly when f\_V(A)=0 together with the already automatic right-invertibility from node 1.1.2; node 1.3.1 identifies this zero locus exactly as A\textasciicircum{}perp=g\_V(A\textasciicircum{}parallel). Finally every U in calU belongs to calUbar\_delta for every delta>0. Taking V=U gives f\_U(0)=0, and uniqueness yields g\_U(0)=0, so U=X\_U(0). Hence these C\textasciicircum{}1 graph charts cover calU. \steptag{archived} \\[4pt]
\step{1.5} Constants-closing step: choose C\_ch=1024 and kappa\_ch=1/4. The guard C\_ch*(epsilon\_r+delta)<=kappa\_ch implies all smallness hypotheses in nodes 1.1-1.4, and their conclusions give every existence, uniqueness, membership, estimate, C\textasciicircum{}1-chart, and covering clause of node 1 with the advertised C\_ch bounds and strict normal-derivative bound <1. \steptag{archived} \\[6pt]
\step{1.5.1} Dependency-gated closure. Set C\_ch=1024 and kappa\_ch=1/4. These are universal, C\_ch>=1, and 0<kappa\_ch<=1/2. For any delta>0 satisfying the root guard, s=epsilon\_r+delta<=1/4096<1/512 and c=64s<=1/64<1/8. Thus nodes 1.1 and 1.2 apply, node 1.3 supplies the unique C\textasciicircum{}1 map g\_V on the whole anti-Hermitian 2delta-ball, and node 1.4 supplies the unitary graph points and covering. The quantitative bounds strengthen to ||g\_V(P)+q/2||<=64(epsilon\_r delta+delta\textasciicircum{}2)<=C\_ch(epsilon\_r delta+delta\textasciicircum{}2), ||Dg\_V(P)||<=128s<=C\_ch s, and ||D\_HF-I||<=64s<=C\_ch s<=1/4<1. Since q=V\textasciicircum{}dagger bold-dot V-J and F(P,H)=f\_V(P+H), these are literally the three estimates and all remaining clauses in node 1. \steptag{archived} \\[4pt]
\end{proofbox}

\end{document}
